\documentclass{article}
\usepackage{hyperref}
\usepackage{datetime}
\usepackage{graphicx}
\usepackage{listings}
\usepackage{color}
\usepackage{todonotes}
\usepackage{lmodern}
\usepackage[T1]{fontenc}
\usepackage{scrextend}
\addtokomafont{labelinglabel}{\sffamily}
\usepackage{blindtext}
\usepackage{marginnote}
\usepackage{hyperref}
\usepackage{tikz}
\usepackage{subfiles}
\usepackage{import}
\usetikzlibrary{automata, positioning, arrows}
\hypersetup{
    colorlinks=true,
    linktoc=all,
    linkcolor=blue,
    urlcolor=blue,
}
\author{Peter Mezei}
\title{Gardenzilla \\ development notes}
\pagestyle{myheadings}
\markright{Gardenzilla\hfill Development Notes \hfill}
\begin{document}
\maketitle
\thispagestyle{empty}
\newpage
\thispagestyle{empty}
\begin{figure}
  \centering
  \includegraphics[width=3cm]{logo.png}
\end{figure}
\begin{center}
  {\textit  Gardenzilla development notes, questions and future development directions. The main goal of this document is to document the development steps and collect all the ideas for future evaluation.

    \vspace{1cm} All right reserved \the\year{}

    All the logos, images and all the other\\ used visual materials are copyrighted.
  }
\end{center}
\newpage
\thispagestyle{empty}
\tableofcontents
\newpage
\section{Introduction}

\subsection{License}

\marginnote{License \\GNU GPL\\ version 2}

We mainly use GNU/GPL version 2. This means this is a free software,
you can freely use it, and modify it. IMPORTANT! As GNU/GPL version 2
is a COPYLEFT license, YOU MUST USE THE SAME GNU/GPL2 LICENSE in your
own verion. If you modify any of the project, YOU MUST MAKE YOUR
UPDATED VERSION AVAILABLE ONLINE. As this project is licensed under
GNU/GPLv2, you most apply the same license and use it - or publish it
- under the same GNU/GPLv2 licesen. If you use this project as a smaller
part of a larger `product`
\footnote{Software, Service, or any kind of
  business outcome where you use this product}
YOU MUST APPLY THE SAME GNU/GPLv2 LICENSE FOR YOUR LARGER PRODUCT AS WELL.

\subsection{Technology}

We use Rust as the main programming language. It is extremely powerfull,
effective, blazangly fast and safe by design. No race condition, and
almost impossible to create any bug that compiles. No garbage collection,
no runtime dependencies.

Other core technolgoies inside the project:

\begin{labeling}{alligator}
  \item [Gnu/Linux] \hfill \\ main development and hosting environment
  \item [Rust] Main programming language
  \item [GnuMakefile] \hfill \\ Using for main build and cleanup staffs
  \item [LaTex] for documentaion
  \item [Nginx] for webserver
  \item [Docker] for container management
  \item [TypeScript] \hfill \\ for typesafe javascript editing
  \item [HTML] for web documents
  \item [CSS] for web styles
  \item [SASS] for better css creation
  \item [JavaScript] \hfill \\ for client development
  \item [Angular] as the main client library
  \item [JSON] as the main API communication format
\end{labeling}

\subsection{3rd party libraries}

There are plenty awesome 3rd party libreries that we heavily use.
We would like to thank you for all the developers from the open-source
community who and whose projects helped us a lot.
In the above you can see the major 3rd party libs we use:

\subsubsection{3rd party Rust libraries}

\begin{labeling}{alligator}
  \item [storaget] data object layer as DAL. In the core project we do
    not use database, instead we use rust data objects inside
    storaget::Pack<T> and storaget::VecPack<T>. Pack manages all the
    data sync between memory and disk.\footnote{Storaget is our inhouse sub-project. See: \url{https://github.com/mezeipetister/storaget}}
  \item [chrono] date library. Using to manage date and datetimes. In
    the core project we use Date<Utc> and DateTime<Utc>. Whenever possible
    we avoid using NaiveDate(s).
  \item [rand] for random number generation.
  \item [bcrypt] for HTTP auth TOKEN encryption.
  \item [lettre] for SMTP email sending.
  \item [nanoid] for ID generation.
  \item [rocket] web framework for REST API.
  \item [rocket\_cors] \hfill \\ cors management for REST API
  \item [serde] object serialization, deserialization
  \item [jwt] for JSON web token management
\end{labeling}

\subsection{3rd party NodeJS libraries}

\begin{labeling}{alligator}
  \item [Angular] core web framework
  \item [chart.js] displaying charts
  \item [ng2-chartjs] \hfill \\ chart.js -> Agnular integration
  \item [markdown.it] \hfill \\ markdown management
  \item [markdown] \hfill \\ markdown management
\end{labeling}

\newpage

% section MODULES
% ***************
\section{Modules}

\subsection{Customers}

We manage customers as independent data objects, and manage a customer
object through its implemented methods, or via public functions inside
core::customer module.

\begin{figure}[!htb]
  \caption{Customer object}
  \begin{verbatim}
    pub struct Customer {
      id: String,
      related_users: Vec<String>,
      name: String,
      tax_number: String,
      address: InvoiceAddress,
      phone: String,
      email: String,
      date_created: DateTime<Utc>,
      created_by: String,
    }
  \end{verbatim}
\end{figure}

Where:

\begin{description}
  \item[id] \hfill
    \\ Customer unique identifier \todo {Describe and set the way we manage IDs.}
  \item[related\_users] \hfill
    \\ vector of related user ID. User and Customer objects are independent
    of each other, as we can have customer without a user, and we can have
    - at least in theory - a user without a customer. Customer can
    purchase, User can login and interact withing the client. To manage
    Customer <-> User relations, we store them as a related ID list in
    each other. Here in a Customer object thats why we have this
    related\_users vector.

  \item[name] \hfill
    \\ name as string of the customer

  \item[tax\_number] \hfill
    \\ tax number as string
  \item[address] \hfill
    \\ address as InvoiceAddress object \todo{display InoviceAddress}
  \item[phone] \hfill
    \\ phone number as string
  \item[email] \hfill
    \\ email as string
  \item[date\_created] \hfill
    \\ date as chrono::DateTime<Utc>
  \item[created\_by] \hfill
    \\ userid as string, who created this object \todo{describe datecreated
      and created by somewhere as we use them everywhere}
\end{description}

\subsubsection{Methods}

\begin{description}
  \item[new(\dots) -> Self] \hfill
    \\ creates a new instance of Customer. ID cannot be changed after
    creation
  \item[get\_id() -> \&str] \hfill
    \\ returns customer id as \&str
  \item[get\_name() -> \&str] \hfill
    \\ returns customer name as \&str
  \item[set\_name(name: String)] \hfill
    \\ set customer name to the given new one
  \item[has\_user() -> bool] \hfill
    \\ returns bool if there is at least one related user
  \item[get\_users() -> \&Vec<String>] \hfill
    \\ returns related user vector reference
  \item[remove\_user(userid: \&str)] \todo{should return bool, or result} \hfill
    \\ remove a user by its userid
  \item[get\_tax\_number() -> \&str] \hfill
    \\ returns tax number as \&str
  \item[set\_tax\_number(tax\_number: String)] \hfill
    \\ set tax number
\end{description}

\newpage
\todototoc
\listoftodos

\newpage
\thispagestyle{empty}
\vspace*{5cm}
\begin{center}
  {\textit Gardenzilla\\ Development Notes \\ \vspace{5mm} All right reserved \\ Peter Mezei \\\the\year{}}
\end{center}
\end{document}
